\documentclass[12pt]{article}
\usepackage[utf8]{inputenc}
\usepackage[T1]{fontenc}
\usepackage{hyperref}
\title{Berawan, Central Language Resources for K-12 Curriculum Development}
\author{Compiled by Claude Sonnet 5.5 (claude-sonnet-5-5) for Andr\'as Kornai}
\date{}
\begin{document}
\maketitle

\section{Introduction}
Central Berawan (ISO 639-3 zbc; Glottocode cent2098) is an endangered Austronesian language of the Baram and Tutoh area of Sarawak, Malaysia. Berawan was formerly a single ISO code (lod), later split into West (zbw), Central (zbc) and East (zbe). The literature describes four Berawan lects: Batu Belah, Long Teru, Long Jegan and Long Terawan. Sources disagree on exactly how these lects map to the three codes and on population figures, so the details below should be checked. Joshua Project lists about 1,900 Central Berawan, at Batu Belah and Long Teru.

\section{Primary Educational Resources}
No school textbooks or primers were found. The most relevant item is a paper on a unified Berawan orthography for the four lects, developed with Berawan communities, which would be a basis for teaching materials.
\begin{itemize}
\item One Orthography, Four Lects: The Unified Berawan Orthography. \url{https://publisher.unimas.my/ojs/index.php/BJK/article/view/2411}
\item Full text of the same article. \url{https://publisher.unimas.my/ojs/index.php/BJK/article/view/2411/1012}
\end{itemize}

\section{Linguistic and Reference Materials}
\begin{itemize}
\item Wikipedia article on the Berawan language. \url{https://en.wikipedia.org/wiki/Berawan_language}
\item Robert Blust recording collection RB1-003, including Berawan material from 1971. \url{https://researchdata.edu.au/rb1-003-robert-blust-recording/1577760}
\item PARADISEC catalogue item RB1/003. \url{https://catalog.paradisec.org.au/repository/RB1/003}
\item Sarawak Museum web page on Berawan. \url{https://museum.sarawak.gov.my/web/subpage/webpage_view/847}
\item Burkhardt, paper on Berawan diphthongs (Terawan Berawan). \url{https://sil-philippines-languages.org/ical/papers/burkhardt-terawan_berawan_dipthong.pdf}
\end{itemize}

\section{Religious and Community Resources}
\begin{itemize}
\item Global Recordings Network, Berawan, Central (Long Teru). \url{https://globalrecordings.net/ja/language/7983}
\item Mirror of the same page. \url{https://glorecnet-website.k.globalrecordings.net/en/language/7983}
\item Joshua Project, Central Berawan in Malaysia. \url{https://joshuaproject.net/people_groups/19812}
\end{itemize}
Joshua Project indicates only Bible portions and one gospel recording exist.

\section{Digital and Audio Resources}
The Blust recordings (see above) are the main audio found, along with the Global Recordings material.

\section{Curriculum Development Context}
The language is spoken as a first language mainly by older adults and is not known to be taught in schools, so a curriculum would be revitalization-oriented. The unified orthography is a starting point. Community consultation is essential. No existing school-level materials were found.
\end{document}